\documentclass[12pt]{article}
\usepackage[utf8]{inputenc}
\usepackage[T1]{fontenc}
\usepackage{hyperref}
\title{Hanga Hundi Language Resources for K-12 Curriculum Development}
\author{Compiled by Claude Sonnet 5.5 (claude-sonnet-5-5) for Andr\'as Kornai}
\date{}
\begin{document}
\maketitle

\section{Introduction}
Hanga Hundi (ISO 639-3: wos; Glottocode: hang1263) is also called Kwasengen. It is a Ndu language of the Sepik family in the Sepik River region of northern Papua New Guinea. It is spoken in the villages of Nungwaia, Dumek, Weko and Warmetali in Pagwi District, according to Omniglot. Estimates are about 7,200 speakers (2000 census) or about 7,000.

\section{Primary Educational Resources}
I found no primers or school materials in search results.

\section{Linguistic and Reference Materials}
\begin{itemize}
\item Thomas Wendel (1993), \emph{A Preliminary Grammar of Hanga Hundi}, master's thesis, University of Texas at Arlington. The reference record is at \url{https://glottolog.org/resource/reference/id/97333}.
\item Wikipedia article (under the name Kwasengen): \url{https://en.wikipedia.org/wiki/Kwasengen_language}
\item Omniglot page on Kwasengen (Hanga Hundi): \url{https://omniglot.com/writing/kwasengen.htm}
\item WALS entry: \url{https://wals.info/languoid/lect/wals_code_hhu}
\item Grambank entry: \url{https://grambank.clld.org/languages/hang1263}
\item ASJP word list: \url{https://asjp.clld.org/languages/HANGA_HUNDI}
\item Numerals page (Max Planck Institute for Evolutionary Anthropology): \url{https://lingweb.eva.mpg.de/channumerals/Hanga-Hundi.htm}
\end{itemize}

\section{Religious and Community Resources}
A Wycliffe New Testament and Jonah translation exists in Hanga Hundi (1997, released 2013 under a Creative Commons license), available as PDFs from eBible.
\begin{itemize}
\item Scripture book PDF: \url{https://ebible.org/pdf/wos/wos_book.pdf}
\item Scripture, complete PDF: \url{https://ebible.org/pdf/wos/wos_all.pdf}
\end{itemize}

\section{Digital and Audio Resources}
No audio or video resources were found in the searches.

\section{Curriculum Development Context}
The language has about 7,000 speakers and a written scripture, which gives an orthographic base. The grammar is a 1993 thesis, and I found no dictionary or school materials. A curriculum would need the orthography confirmed, a lexicon and graded reading materials, probably built on the scripture text.
\end{document}